\section{Introduction}

As deep neural network architectures scale in size, their representation power grows. However, this capacity expansion demands massive compute and memory during training and inference \citep{hoffmann2022training}. Network pruning has therefore emerged as a fundamental subject, both for practical model deployment \citep{lecun1989optimal, hassibi1993second, reed1993pruning, han2015learning, sze2017efficient} and for gaining theoretical insights into artificial generalisation bounds \citep{arora2018stronger} and biological neural computation \citep{tanaka2019from}.

Existing pruning techniques generally fall into three categories. Post-hoc methods train a full model first, prune low-importance weights, and fine-tune the rest. This requires expensive prune-retrain cycles \citep{han2016deep}. Static methods prune connections once at initialisation before training. But their accuracy remains strictly bounded by initial weight sensitivity \citep{lee2019snip, tanaka2020pruning}. Dynamic sparse training methods prune and regrow weights during training \citep{evci2020rigging}. But they force rigid, pre-specified sparsity targets onto each layer, preventing the model from finding its natural capacity.

In contrast to artificial networks, biological brains achieve extraordinary computational efficiency through dynamic, highly sparse connectivity. A driver of biological learning is structural plasticity via activity-dependent synapse elimination \citep{hebb1949organization, han2024developmental}. Neurobiological studies show that active pathways are preserved while inactive connections are eliminated \citep{faust2021mechanisms, mikuni2013arc, yasuda2021activity}. During early development, the brain grows excess synapses and prunes those that stay idle or fail to support functional circuits \citep{huttenlocher1979synaptic, huttenlocher2002neural}. Inspired by this process, we propose \textbf{Dynamic Activity-Dependent Pruning (DADP)}, a structural plasticity algorithm grounded in a reverse Hebbian mechanism. Rather than evaluating weights by magnitude alone, DADP continuously measures connection importance during training as the expected product of pre-synaptic activation and post-synaptic error gradient: $I_{ij} = \mathbb{E}\left[ \left| a_i \cdot \frac{\partial L}{\partial y_j} \right| \right]$. A connection survives only if its pre-synaptic neuron fires ($a_i \neq 0$) AND its post-synaptic target is actively reducing error ($\frac{\partial L}{\partial y_j} \neq 0$).

Our primary contributions are summarized as follows:
\begin{itemize}
    \item \textbf{Reverse Hebbian Pruning Metric}: We introduce a dynamic structural plasticity criterion ($I_{ij} = \mathbb{E}[|a_i \cdot \frac{\partial L}{\partial y_j}|]$) that evaluates synaptic importance via pre-synaptic activation and backpropagated loss sensitivity within a single-pass training framework.
    \item \textbf{Self-Regulating Sparsity Equilibrium}: We show that DADP operates under a single global threshold $\tau$ without manual layer targets. This establishes a self-correcting negative feedback loop (prune-recovery cycle) that organically discovers optimal layer-wise capacity allocations.
    \item \textbf{Emergent Structural \& Topological Plasticity}: We demonstrate that unstructured connection pruning naturally induces structured filter/neuron death while automatically protecting critical residual skip-connections and maintaining continuous Gaussian weight distributions (avoiding the artificial bimodal exclusion gap of magnitude pruning).
    \item \textbf{Initialization Scheme Invariance}: We establish that DADP exhibits tight structural and performance invariance across standard variance-preserving initialization protocols (Kaiming Normal/Uniform, Xavier Normal/Uniform, Orthogonal).
    \item \textbf{Pruning Interval ($\Delta t$) Robustness}: We demonstrate hyperparameter stability across a 50$\times$ range of pruning intervals ($\Delta t \in [100, 5000]$ batches), proving DADP is not hyperparameter-sensitive.
    \item \textbf{Extensive Benchmarking \& Representation Quality}: Across vision (VGG-16, ResNet-18), NLP/sequence paradigms (BiLSTM-CRF, MiniBERT), and MLPs, DADP matches or exceeds baseline algorithms (Magnitude, SNIP, RigL) at extreme sparsities ($>99\%$ on ResNet-18) while preserving matrix-based dataset entropy and effective rank.
\end{itemize}
